%! Unit 1, Lessons 1.0–1.7 — SAMPLE Test (blank)
% The practice form of the Unit 1 test. Item for item it is the parallel form of ../actual_test/:
% the same 15 multiple-choice targets in the same order and the same five free-response
% questions with the same parts and points (MC 15 x 4 = 60, FRQ 5 x 8 = 40, total 100;
% 20 questions; 5 pages at most) --- every context different, so working this one teaches the
% moves, not the answers.
% Merged into the unit student packet via ../../sample_test/ (see shared/tests.mk).
% NAME ROW: \namedateperiod IS carried --- a unit test is taken in a testing setting.
% NO SKETCH ITEMS: every display is pre-drawn in TikZ and read.
% COMPACT LAYOUT: short answer choices run in two columns (multicols), figures sit beside their
% stems, and each free-response part gets only the lines its answer needs. Keep it to 5 pages.
%
% BLUEPRINT (the free response carries individuals, classifying variables, statistical
% questions, relative frequency, and the 1.5 x IQR rule, so the multiple choice does not repeat them)
%   MC 1 grouping (binning)               gym workout minutes              1.1
%   MC 2 groups of unequal size           two summer camps, swimming       1.2
%   MC 3 discrete vs continuous           mile time to the second          1.3
%   MC 4 read a stem plot (median)        15 April high temperatures       1.3
%   MC 5 interval width                   50 battery lifetimes             1.3
%   MC 6 shape from a histogram           40 used-textbook prices          1.4
%   MC 7 unusual features (dotplot)       20 song lengths                  1.4
%   MC 8 mean vs median from shape        80 home sale prices              1.5
%   MC 9 resistance (a typo)              9 lunch-special prices           1.5
%   MC10 standard deviation in context    30 dog-food bags                 1.5
%   MC11 percentile                       state reading test               1.5
%   MC12 every section holds a quarter    30 basketball games              1.6
%   MC13 parallel boxplots                two urgent-care clinics          1.6
%   MC14 parameter vs statistic           town residents' ages             1.7
%   MC15 z-scores across distributions    history vs chemistry tests       1.7
%   FRQ1 individuals, variables, SQ       city playgrounds                 1.0, 1.1, 1.3
%   FRQ2 relative frequency comparison    two high schools' senior plans   1.2
%   FRQ3 histogram, median, describe           50 trout lengths                 1.3, 1.4
%   FRQ4 five numbers, outlier, center    15 used video-game prices        1.5, 1.6
%   FRQ5 empirical rule, z-scores         egg weights                      1.7
% ANSWERS (MC): 1E 2A 3D 4B 5E 6C 7A 8D 9C 10E 11A 12B 13D 14C 15B
%
% THE DATA (every number checked)
%  MC2  Camp A 54/120 = 0.45; Camp B 39/75 = 0.52.
%  MC4  54 56 58 59 60 61 63 63 65 67 68 70 72 72 76; median (8th) 63; mean 964/15 = 64.3.
%  MC6  bins of $10 from 0: 12 11 7 5 3 1 1 (40), skewed right.
%  MC7  2.5:1 3.0:3 3.5:6 4.0:5 4.5:3 5.0:1 8.0:1 (20).
%  MC9  6 7 7 8 9 9 10 11 12 -> 120: Q1 7, med 9, Q3 10.5 both ways; mean, range move.
%  MC12 52/60/64/75/88.
%  MC13 A 10/18/24/30/45 (IQR 12); B 15/26/30/34/40 (IQR 8).
%  MC15 Maya (84-76)/5 = 1.6; Jordan (88-80)/8 = 1.0.
%  FRQ2 North 120/55/50/25 (250) = 48/22/20/10%; South 88/32/28/12 (160) = 55/20/17.5/7.5%.
%  FRQ3 bins of 4 cm from 12: 2 0 3 5 9 15 11 5 (50); cum 2 2 5 10 19 34 45 50.
%  FRQ4 8 10 12 12 14 15 15 16 18 20 21 22 24 25 58: Q1 12, med 16, Q3 22, IQR 10,
%       fences -3 and 37; sum 290, mean 19.33.
%  FRQ5 N(58, 4): 50-66 95%; >62 16%; 51 -> z -1.75; 65 -> 1.75 vs farm 2 N(64, 6) 73 -> 1.5.
%
% LOCKSTEP: GENERATED from its answer key under ../../test_keys/ --- edit the key, then regenerate.
% Every ans mark on an MC option prints bare there; every other ans mark becomes a phantom of
% itself; every run of n anslines becomes writelines{n}; work blocks stay byte-identical.
\documentclass[12pt]{article}
\usepackage{apstats-article}
\usepackage{apstats-boxes}
\usepackage{multicol}
\setlength{\multicolsep}{1pt}

% Answer choices: tight list; wrap in multicols{2} when every choice fits half a line.
\newenvironment{choices}
  {\begin{enumerate}[label=\textbf{(\Alph*)}, leftmargin=2.2em, itemsep=0pt, topsep=1pt]}
  {\end{enumerate}}

\begin{document}

\pageheader{Unit 1 \quad Lessons 1.0--1.7}{Sample Test}

\namedateperiod
\par\vspace{2pt}
\hfill\small Score:\;\blank{1.3cm}\;/\;100\normalsize
\vspace{0.04in}

\boxguard[6]
\begin{headlinebox}{goldbg}
\small
\textbf{Section I --- Multiple Choice} \hfill \textbf{60 points, 4 each}\\
This is the \textbf{practice} test --- not scored. Circle the single best answer.
\end{headlinebox}

\vspace{0.06in}

\begin{enumerate}[leftmargin=1.9em, itemsep=0.03in]

  \boxguard[8]
  \item A gym records each member's weekly workout time in minutes, then replaces it with a
        label: \emph{Light} (under 90), \emph{Moderate} (90 to 180), or \emph{Heavy} (over 180).
        Which statement is true?
        \begin{choices}
          \item The new variable is still quantitative, because the labels are based on minutes.
          \item The new variable is quantitative, because the labels are in order.
          \item Nothing is lost: the minutes can be recovered from the labels.
          \item The median workout time can still be found exactly from the labels.
          \item It is now categorical, and the mean workout time can no longer be found.
        \end{choices}

  \boxguard[8]
  \item At Camp Alder, 54 of 120 campers named swimming their favorite activity. At Camp Birch,
        39 of 75 did. Which is best supported?
        \begin{multicols}{2}
        \begin{choices}
          \item Birch: 52\% is greater than 45\%
          \item Alder: 54 is greater than 39
          \item Equally popular at both camps
          \item Can't compare: camps differ in size
          \item Alder: it has more campers
        \end{choices}
        \end{multicols}

  \boxguard[6]
  \item Which of these variables is \textbf{continuous}?
        \begin{multicols}{2}
        \begin{choices}
          \item Number of siblings
          \item Number of texts sent today
          \item Number of pets at home
          \item Mile time, to the nearest second
          \item Number of songs on a playlist
        \end{choices}
        \end{multicols}

  \boxguard[8]
  \item \begin{minipage}[t]{0.62\linewidth}\raggedright
        The stem plot shows the daily high temperature ($^\circ$F) in a city on 15 days in
        April. What is the \textbf{median} daily high?
        \begin{multicols}{3}
        \begin{choices}
          \item $61^\circ$F
          \item $63^\circ$F
          \item $64.3^\circ$F
          \item $65^\circ$F
          \item $7^\circ$F
        \end{choices}
        \end{multicols}
        \end{minipage}\hfill
        \begin{minipage}[t]{0.32\linewidth}
        \centering\small
        \begin{tabular}[t]{@{} r | l @{}}
        5 & 4 6 8 9 \\
        6 & 0 1 3 3 5 7 8 \\
        7 & 0 2 2 6 \\
        \end{tabular}\par\vspace{2pt}
        {\footnotesize Key: $5\,|\,8 = 58^\circ$F}
        \end{minipage}

  \boxguard[8]
  \item Two histograms of the \emph{same} 50 battery lifetimes (hours): with 1-hour intervals there
        is a gap from 12 to 14 hours; with 4-hour intervals there is no gap. Which is correct?
        \begin{choices}
          \item One of the two histograms must have been drawn incorrectly.
          \item The wider intervals changed some of the battery lifetimes.
          \item Only the histogram with more bars is a true histogram.
          \item The gap shows that the data set contains an outlier.
          \item Both are correct: the interval width changes the picture, not the data.
        \end{choices}

  \boxguard[9]
  \item \begin{minipage}[t]{0.56\linewidth}\raggedright
        The histogram shows the prices of 40 used textbooks sold online. Which best describes
        the \textbf{shape}?
        \begin{choices}
          \item Skewed left: most prices pile up on the left.
          \item Roughly symmetric: there is one peak.
          \item Skewed right: a few prices are high.
          \item Uniform: every interval has a textbook.
          \item Bimodal: the two lowest bars are tall.
        \end{choices}
        \end{minipage}\hfill
        \begin{minipage}[t]{0.40\linewidth}
        \centering
        \begin{tikzpicture}[x=0.062cm, y=0.19cm, baseline=(current bounding box.north)]
          \foreach \y in {0,4,8,12} {
            \draw[linegray, very thin] (0,\y) -- (72,\y);
            \node[left, font=\tiny] at (0,\y) {\y};}
          \foreach \l/\c in {0/12, 10/11, 20/7, 30/5, 40/3, 50/1, 60/1}
            \draw[fill=skymid] (\l,0) rectangle (\l+10,\c);
          \draw[thick] (0,0) -- (73,0);
          \draw[thick] (0,0) -- (0,13);
          \foreach \x in {0,10,...,70} { \node[below, font=\tiny] at (\x,0) {\x}; }
          \node[font=\tiny] at (36,-3.2) {Price (dollars)};
          \node[rotate=90, font=\tiny] at (-10,6.5) {Count};
        \end{tikzpicture}
        \end{minipage}

  \boxguard[9]
  \item \begin{minipage}[t]{0.56\linewidth}\raggedright
        The dotplot shows the lengths (minutes) of the 20 songs on a playlist. Which statement
        is best?
        \end{minipage}\hfill
        \begin{minipage}[t]{0.40\linewidth}
        \centering
        \begin{tikzpicture}[x=0.95cm, y=0.2cm, baseline=(current bounding box.north)]
          \draw[thick] (2,0) -- (8.5,0);
          \foreach \x in {2,3,...,8} {
            \draw (\x,-0.4) -- (\x,0.4); \node[below, font=\tiny] at (\x,-0.3) {\x};}
          \foreach \x/\n in {2.5/1, 3/3, 3.5/6, 4/5, 4.5/3, 5/1, 8/1}
            \foreach \k in {1,...,\n} \fill[navy] (\x,\k*0.9) circle (2pt);
          \node[font=\tiny] at (5.25,-3) {Song length (minutes)};
        \end{tikzpicture}
        \end{minipage}
        \begin{choices}
          \item The 8-minute song is a possible outlier, set off by a gap; investigate it.
          \item The 8-minute song is an outlier, so it should be deleted from the data.
          \item There are no unusual features, because most songs last 3 to 4.5 minutes.
          \item The distribution is bimodal, with peaks at 3.5 and 8 minutes.
          \item The gap from 2.5 to 3 minutes is the most unusual feature.
        \end{choices}

  \boxguard[8]
  \item The distribution of 80 home sale prices in a town is strongly \textbf{skewed right}.
        Which is most likely true?
        \begin{multicols}{2}
        \begin{choices}
          \item Mean $<$ median; report the mean
          \item Mean $=$ median
          \item Mean $>$ median; report the mean
          \item Mean $>$ median; report the median
          \item Can't tell without the data
        \end{choices}
        \end{multicols}

  \boxguard[8]
  \item The prices of 9 lunch specials are \$6, 7, 7, 8, 9, 9, 10, 11, and 12. The \$12 price is
        mistakenly entered as \$120. Which summaries change?
        \begin{multicols}{2}
        \begin{choices}
          \item Both the mean and the median
          \item The median, but not the mean
          \item Mean and range; not median or IQR
          \item None of them: only one price is wrong
          \item Only the IQR
        \end{choices}
        \end{multicols}

  \boxguard[8]
  \item For 30 bags of dog food labeled 20 lb, the standard deviation of the weights is 0.3 lb.
        Which is the best interpretation?
        \begin{choices}
          \item Every bag weighs within 0.3 lb of the mean weight.
          \item A typical bag weighs 0.3 lb.
          \item The heaviest bag weighs 0.3 lb more than the lightest bag.
          \item About 30\% of the bags weigh more than the mean weight.
          \item The bag weights typically vary by about 0.3 lb from the mean weight.
        \end{choices}

  \boxguard[8]
  \item A student's score on a state reading test is at the \textbf{82nd percentile}. This means
        \begin{choices}
          \item about 82\% of the test-takers scored at or below this student.
          \item the student answered 82\% of the questions correctly.
          \item the student scored 82 points.
          \item about 18\% of the test-takers scored at or below this student.
          \item the student scored higher than exactly 82 other students.
        \end{choices}

  \boxguard[9]
  \item \begin{minipage}[t]{0.56\linewidth}\raggedright
        The boxplot shows the points a basketball team scored in each of its 30 games. Which
        statement is true?
        \end{minipage}\hfill
        \begin{minipage}[t]{0.40\linewidth}
        \centering
        \begin{tikzpicture}[x=0.15cm, y=1cm, baseline=(current bounding box.north)]
          \draw[thick] (50,0) -- (90,0);
          \foreach \x in {50,60,...,90} {
            \draw (\x,-0.07) -- (\x,0.07); \node[below, font=\tiny] at (\x,-0.05) {\x};}
          \draw (52,0.45) -- (60,0.45);
          \draw (75,0.45) -- (88,0.45);
          \draw[fill=skymid] (60,0.25) rectangle (75,0.65);
          \draw[thick] (64,0.25) -- (64,0.65);
          \draw (52,0.33) -- (52,0.57);
          \draw (88,0.33) -- (88,0.57);
          \node[font=\tiny] at (70,-0.45) {Points scored};
        \end{tikzpicture}
        \end{minipage}
        \begin{choices}
          \item More games scored 64 to 75 than 60 to 64, because that part of the box is wider.
          \item About a quarter of the games scored 60 to 64, and about a quarter 64 to 75.
          \item About half of the games had scores above 75.
          \item The team's mean score was 64 points.
          \item The range of the scores is 28 points.
        \end{choices}

  \boxguard[10]
  \item \begin{minipage}[t]{0.50\linewidth}\raggedright
        The parallel boxplots show patients' waiting times (minutes) at two urgent-care clinics.
        Which comparison is best?
        \end{minipage}\hfill
        \begin{minipage}[t]{0.46\linewidth}
        \centering
        \begin{tikzpicture}[x=0.14cm, y=1cm, baseline=(current bounding box.north)]
          \draw[thick] (5,0) -- (50,0);
          \foreach \x in {10,20,...,50} {
            \draw (\x,-0.07) -- (\x,0.07); \node[below, font=\tiny] at (\x,-0.05) {\x};}
          \draw (10,1.0) -- (18,1.0); \draw (30,1.0) -- (45,1.0);
          \draw[fill=skymid] (18,0.82) rectangle (30,1.18); \draw[thick] (24,0.82) -- (24,1.18);
          \draw (10,0.9) -- (10,1.1); \draw (45,0.9) -- (45,1.1);
          \node[left, font=\tiny] at (5,1.0) {A};
          \draw (15,0.45) -- (26,0.45); \draw (34,0.45) -- (40,0.45);
          \draw[fill=goldbg] (26,0.27) rectangle (34,0.63); \draw[thick] (30,0.27) -- (30,0.63);
          \draw (15,0.35) -- (15,0.55); \draw (40,0.35) -- (40,0.55);
          \node[left, font=\tiny] at (5,0.45) {B};
          \node[font=\tiny] at (27.5,-0.45) {Waiting time (minutes)};
        \end{tikzpicture}
        \end{minipage}
        \begin{choices}
          \item Clinic A has the higher median wait, because its box reaches farther right.
          \item The clinics are equally variable, because the boxes are the same height.
          \item Clinic A's waits are longer, because its maximum is larger.
          \item B has the higher median (30 vs.\ 24 min) and the smaller IQR (8 vs.\ 12).
          \item Every patient at Clinic B waited longer than every patient at Clinic A.
        \end{choices}

  \boxguard[8]
  \item A town's records show the mean age of all 4,200 residents is 38.6 years. A reporter
        surveys 50 residents; their mean age is 41.2 years. Which is correct?
        \begin{multicols}{2}
        \begin{choices}
          \item 38.6 is $\bar x$; 41.2 is $\mu$
          \item Both are parameters
          \item 38.6 is $\mu$; 41.2 is $\bar x$
          \item Both are statistics, since both are means
          \item 41.2 is $\mu$, since it is larger
        \end{choices}
        \end{multicols}

  \boxguard[9]
  \item Maya scored 84 on a history test (class mean 76, standard deviation 5). Jordan scored 88 on
        a chemistry test (class mean 80, standard deviation 8). Who did better \emph{relative to
        their own class}?
        \begin{multicols}{2}
        \begin{choices}
          \item Jordan: 88 is greater than 84
          \item Maya: $z = 1.6$ vs.\ Jordan's 1.0
          \item Equal: both are 8 points above
          \item Jordan: $z = 1.6$ vs.\ Maya's 1.0
          \item Can't compare different subjects
        \end{choices}
        \end{multicols}

\end{enumerate}

\vspace{0.08in}

\boxguard[12]
\begin{headlinebox}{goldbg}
\small
\textbf{Section II --- Free Response} \hfill \textbf{40 points, 8 each}\\
Show your reasoning. Answer \emph{in context} --- name the variables and their units.
\end{headlinebox}

\vspace{0.06in}

\begin{enumerate}[leftmargin=1.9em, itemsep=0.10in, resume]

  \boxguard[8]
  \item \textbf{(8 points)} A city's parks department keeps these records on its playgrounds.

        \vspace{3pt}
        {\centering\small
        \renewcommand{\arraystretch}{1.05}
        \begin{tabular}{@{} c l c c l @{}}
        \textbf{Park \#} & \textbf{Neighborhood} & \textbf{Swings} & \textbf{Area (sq ft)} &
        \textbf{Surface} \\ \hline
        104 & Eastgate   & 6 & 2,450 & Wood chips \\
        117 & Mill Creek & 4 & 1,820 & Rubber \\
        122 & Eastgate   & 8 & 3,100 & Rubber \\
        135 & Riverside  & 2 & 960   & Sand \\
        \end{tabular}\par}

        \vspace{3pt}
        \textbf{(a)} What are the individuals? Is \emph{Eastgate} a variable or a value?\\[2pt]
        \writelines{1}

        \vspace{3pt}
        \textbf{(b)} Classify \emph{Park \#}, \emph{Swings}, and \emph{Area} as categorical or
        quantitative. For each quantitative one, say discrete or continuous.\\[2pt]
        \writelines{2}

        \vspace{3pt}
        \textbf{(c)} Is ``How many swings does Park 122 have?'' a statistical question? Explain,
        then write one that is.\\[2pt]
        \writelines{2}

  \boxguard[14]
  \item \textbf{(8 points)} Two high schools asked their seniors about their plans after
        graduation.

        \vspace{3pt}
        \begin{minipage}[c]{0.46\linewidth}
        \centering\small
        \setlength{\tabcolsep}{4pt}\renewcommand{\arraystretch}{1.05}
        \begin{tabular}{@{} l | c c @{}}
         & \textbf{North High} & \textbf{South High} \\ \hline
        4-year college & 120 & 88 \\
        2-year college & 55  & 32 \\
        Work           & 50  & 28 \\
        Military       & 25  & 12 \\ \hline
        \textbf{Total} & 250 & 160 \\
        \end{tabular}
        \end{minipage}\hfill
        \begin{minipage}[c]{0.50\linewidth}
        \centering
        \begin{tikzpicture}[x=0.62cm, y=0.017cm]
          \foreach \y in {0,40,80,120} {
            \draw[linegray, very thin] (0,\y) -- (8.4,\y);
            \node[left, font=\tiny] at (0,\y) {\y};}
          \foreach \i/\n/\s in {0/120/88, 1/55/32, 2/50/28, 3/25/12} {
            \fill[navy]    (0.3+2*\i,0) rectangle (1.0+2*\i,\n);
            \fill[goldacc] (1.0+2*\i,0) rectangle (1.7+2*\i,\s);}
          \draw[thick] (0,0) -- (8.4,0);
          \draw[thick] (0,0) -- (0,130);
          \foreach \i/\lab in {0/4-yr, 1/2-yr, 2/Work, 3/Military}
            \node[below, font=\tiny] at (1.0+2*\i,0) {\lab};
          \node[rotate=90, font=\tiny] at (-1.1,65) {Seniors};
          \fill[navy] (4.6,128) rectangle (4.9,138); \node[right, font=\tiny] at (4.9,133) {North};
          \fill[goldacc] (6.5,128) rectangle (6.8,138); \node[right, font=\tiny] at (6.8,133) {South};
        \end{tikzpicture}
        \end{minipage}

        \vspace{3pt}
        \textbf{(a)} What proportion of each school's seniors plan to attend a 4-year college?\\[2pt]
        \writelines{1}

        \vspace{3pt}
        \textbf{(b)} A reporter wrote: ``North sends more of its seniors to 4-year colleges ---
        look how much taller its bar is.'' What is wrong with this reasoning?\\[2pt]
        \writelines{2}

        \vspace{3pt}
        \textbf{(c)} Do \emph{most} North seniors plan on a 4-year college? Most South seniors?
        Justify.\\[2pt]
        \writelines{2}

  \boxguard[14]
  \item \textbf{(8 points)} A biologist measured the lengths (cm) of 50 trout caught in a lake.

        \vspace{3pt}
        \begin{minipage}[t]{0.42\linewidth}
        \centering
        \begin{tikzpicture}[x=0.165cm, y=0.15cm, baseline=(current bounding box.north)]
          \foreach \y in {0,5,10,15} {
            \draw[linegray, very thin] (10,\y) -- (46,\y);
            \node[left, font=\tiny] at (10,\y) {\y};}
          \foreach \l/\c in {12/2, 16/0, 20/3, 24/5, 28/9, 32/15, 36/11, 40/5}
            \draw[fill=skymid] (\l,0) rectangle (\l+4,\c);
          \draw[thick] (10,0) -- (46,0);
          \draw[thick] (10,0) -- (10,16);
          \foreach \x in {12,20,...,44} { \node[below, font=\tiny] at (\x,0) {\x}; }
          \node[font=\tiny] at (28,-3.6) {Length (cm)};
          \node[rotate=90, font=\tiny] at (6.6,8) {Trout};
        \end{tikzpicture}
        \end{minipage}\hfill
        \begin{minipage}[t]{0.55\linewidth}\raggedright
        \textbf{(a)} How many trout were shorter than 28 cm? What percent of the catch is that?\\[2pt]
        \writelines{1}

        \vspace{3pt}
        \textbf{(b)} Can the exact median be found from the histogram? If not, which interval
        holds it? Explain.\\[2pt]
        \writelines{3}
        \end{minipage}

        \vspace{3pt}
        \textbf{(c)} Describe the distribution of trout lengths.\\[2pt]
        \writelines{3}

  \boxguard[14]
  \item \textbf{(8 points)} A resale shop's prices (dollars) for 15 used video games, in order:\\[2pt]
        \hspace*{\fill}8 \ \ 10 \ \ 12 \ \ 12 \ \ 14 \ \ 15 \ \ 15 \ \ 16 \ \ 18 \ \ 20 \ \ 21
        \ \ 22 \ \ 24 \ \ 25 \ \ 58\hspace*{\fill}

        \vspace{2pt}

        \textbf{(a)} Find the five-number summary.
        \begin{work}
          \text{Min} = 8, \quad Q_1 = 12, \quad \text{Median} = 16, \quad Q_3 = 22, \quad
          \text{Max} = 58
        \end{work}

        \textbf{(b)} Use the $1.5 \times \text{IQR}$ rule to decide whether \$58 is an outlier. If
        it is, where does the upper whisker of a modified boxplot end?
        \begin{work}
          \text{IQR} = 22 - 12 = 10, \qquad Q_3 + 1.5(\text{IQR}) = 22 + 15 = 37
        \end{work}
        \writelines{1}

        \vspace{3pt}
        \textbf{(c)} The mean price is \$19.33. Which better describes a typical price, the mean or
        the median? Explain.\\[2pt]
        \writelines{2}

  \boxguard[14]
  \item \textbf{(8 points)} The weights of eggs from a farm are approximately normal, with mean
        $\mu = 58$ grams and standard deviation $\sigma = 4$ grams.

        \vspace{2pt}
        \begin{minipage}[t]{0.40\linewidth}
        \centering
        \begin{tikzpicture}[x=0.48cm, y=1.5cm, baseline=(current bounding box.north)]
          \draw[thick, navy, domain=-3.6:3.6, samples=80, smooth]
            plot (\x, {exp(-\x*\x/2)});
          \draw[thick] (-3.8,0) -- (3.8,0);
          \foreach \k/\lab in {-3/46, -2/50, -1/54, 0/58, 1/62, 2/66, 3/70} {
            \draw (\k,-0.05) -- (\k,0.05); \node[below, font=\tiny] at (\k,-0.03) {\lab};}
          \node[font=\tiny] at (0,-0.38) {Egg weight (grams)};
        \end{tikzpicture}
        \end{minipage}\hfill
        \begin{minipage}[t]{0.57\linewidth}\raggedright
        \textbf{(a)} About what percent of the eggs weigh between 50 and 66 grams?\\[2pt]
        \writelines{1}

        \vspace{3pt}
        \textbf{(b)} About what percent weigh \emph{more than} 62 grams?\\[2pt]
        \writelines{1}
        \end{minipage}

        \vspace{3pt}
        \textbf{(c)} One egg weighs 51 grams. Find its $z$-score and interpret it in context.
        \begin{work}
          z = \frac{51 - 58}{4} = -1.75
        \end{work}
        \writelines{1}

        \vspace{3pt}
        \textbf{(d)} A second farm's eggs have mean 64 g and SD 6 g. Which is heavier \emph{relative
        to its own farm}: a 65-g egg from the first farm or a 73-g egg from the second?
        \begin{work}
          z_{\text{first}} = \frac{65 - 58}{4} = 1.75, \qquad
          z_{\text{second}} = \frac{73 - 64}{6} = 1.5
        \end{work}
        \writelines{1}

\end{enumerate}

\end{document}
